\documentclass[11pt]{article}
\usepackage[utf8]{inputenc}
\usepackage[T1]{fontenc}
\usepackage{amsmath, amssymb, amsfonts, amsthm}
\usepackage{mathtools}
\usepackage{geometry}
\usepackage{setspace}
\usepackage{graphicx}
\usepackage{booktabs}
\usepackage{array}
\usepackage{ragged2e}
\usepackage[colorlinks=true, linkcolor=blue, citecolor=blue, urlcolor=blue]{hyperref}
\usepackage[style=apa,natbib=true]{biblatex}
\addbibresource{references.bib} 

\geometry{a4paper, margin=1in}
\linespread{1.15}

\theoremstyle{plain}
\newtheorem{theorem}{Theorem}[section]
\newtheorem{proposition}[theorem]{Proposition}
\newtheorem{lemma}[theorem]{Lemma}
\newtheorem{corollary}[theorem]{Corollary}

\theoremstyle{definition}
\newtheorem{definition}[theorem]{Definition}
\newtheorem{remark}[theorem]{Remark}

\begin{document}

\section{Base Neuronal Dynamics} \label{sec:interspike}

The presynaptic input signals decay as
\[
    \frac{dI_s}{dt} = -\frac{I_s}{\tau_s}, \qquad
    \frac{dx_{\mathrm{pre}}}{dt} = -\frac{x_{\mathrm{pre}}}{\tau_+}, \qquad
    \frac{dr_{\mathrm{pre}}}{dt} = -\frac{r_{\mathrm{pre}}}{\tau_r}.
\]
These are fully determined by initial conditions set at the last presynaptic spike and enter the postsynaptic system as forcing functions. The postsynaptic dynamics are:
\begin{align}
    \tau_m \frac{dV}{dt} &= -(V - E_L) + R_m\, w\, I_s(t), \label{eq:inter_V} \\[6pt]
    \frac{dy_{\mathrm{post}}}{dt} &= -\frac{y_{\mathrm{post}}}{\tau_-}, \label{eq:inter_ypost} \\[6pt]
    \frac{dE}{dt} &= -\frac{E}{\tau_e}, \label{eq:inter_E} \\[6pt]
    \frac{dr_{\mathrm{post}}}{dt} &= -\frac{r_{\mathrm{post}}}{\tau_r}, \label{eq:inter_rpost} \\[6pt]
    \tau_{\bar{R}} \frac{d\bar{R}}{dt} &= -\bar{R} - \bigl(r_{\mathrm{post}} - \alpha\, r_{\mathrm{pre}}\bigr)^2, \label{eq:inter_Rbar} \\[6pt]
    \frac{dw}{dt} &= \Bigl[-\bigl(r_{\mathrm{post}} - \alpha\, r_{\mathrm{pre}}\bigr)^2 - \bar{R}\Bigr]\, E. \label{eq:inter_w}
\end{align}

The postsynaptic state vector is $\mathbf{x}(t) = \bigl(V,\; y_{\mathrm{post}},\; E,\; r_{\mathrm{post}},\; \bar{R},\; w\bigr) \in \mathbb{R}^6$, and the inputs $\bigl(I_s(t),\; x_{\mathrm{pre}}(t),\; r_{\mathrm{pre}}(t)\bigr)$.

\section{Pre-synaptic spike event} \label{sec:pre-spike}

At each presynaptic spike time $t_{\mathrm{pre}}^k$ (an exogenous input event), the input subsystem variables receive discrete increments:
\[
    I_s \;\leftarrow\; I_s + 1, \qquad
    x_{\mathrm{pre}} \;\leftarrow\; x_{\mathrm{pre}} + 1, \qquad
    r_{\mathrm{pre}} \;\leftarrow\; r_{\mathrm{pre}} + 1.
\]
The eligibility trace acquires an LTD contribution, weighted by the postsynaptic trace at the moment of the spike:
\[
    E \;\leftarrow\; E - \eta_-\, w\, y_{\mathrm{post}}(t_{\mathrm{pre}}^{k-}).
\]
All remaining postsynaptic variables ($V$, $y_{\mathrm{post}}$, $r_{\mathrm{post}}$, $\bar{R}$, $w$) are continuous through $t_{\mathrm{pre}}^k$; inter-spike dynamics \eqref{eq:inter_V}--\eqref{eq:inter_w} resume with the updated $I_s$.

\section{Post-synaptic spike event} \label{sec:post-spike}

At each postsynaptic spike time $t_{\mathrm{post}}^k$, defined by the threshold crossing $V = \theta$ with $dV/dt > 0$ (an endogenous event), the membrane potential resets and the postsynaptic trace variables receive discrete increments:
\[
    V \;\leftarrow\; V_{\mathrm{reset}}, \qquad
    y_{\mathrm{post}} \;\leftarrow\; y_{\mathrm{post}} + 1, \qquad
    r_{\mathrm{post}} \;\leftarrow\; r_{\mathrm{post}} + 1.
\]
The eligibility trace acquires an LTP contribution, weighted by the presynaptic trace at the moment of the spike:
\[
    E \;\leftarrow\; E + \eta_+\,(w_{\mathrm{max}} - w)\, x_{\mathrm{pre}}(t_{\mathrm{post}}^{k-}).
\]
All remaining variables ($\bar{R}$, $w$) are continuous through $t_{\mathrm{post}}^k$. Following the spike, voltage integration is suspended for a refractory period $\tau_{\mathrm{ref}}$, after which it resumes from $V_{\mathrm{reset}}$.
\section{Complete Dynamics} \label{sec:reduced_system}

\begin{align}
    \tau_m \frac{dV}{dt} &= -(V - E_L) + R_m\, w\, I_s(t), \label{eq:reduced_V} \\[6pt]
    \tau_- \frac{dy_{\mathrm{post}}}{dt} &= -y_{\mathrm{post}} + \rho_{\mathrm{post}}(t), \label{eq:reduced_ypost} \\[6pt]
    \tau_e \frac{dE}{dt} &= -E + \eta_+(w_{\mathrm{max}} - w)\, x_{\mathrm{pre}}\, \rho_{\mathrm{post}}(t) - \eta_-\, w\, y_{\mathrm{post}}\, \rho_{\mathrm{pre}}(t), \label{eq:reduced_E} \\[6pt]
    \tau_r \frac{dr_{\mathrm{post}}}{dt} &= -r_{\mathrm{post}} + \rho_{\mathrm{post}}(t), \label{eq:reduced_rpost} \\[6pt]
    \tau_{\bar{R}} \frac{d\bar{R}}{dt} &= -\bar{R} - \bigl(r_{\mathrm{post}} - \alpha\, r_{\mathrm{pre}}\bigr)^2, \label{eq:reduced_Rbar} \\[6pt]
    \frac{dw}{dt} &= \Bigl[-\bigl(r_{\mathrm{post}} - \alpha\, r_{\mathrm{pre}}\bigr)^2 - \bar{R}\Bigr]\, E, \label{eq:reduced_w}
\end{align}

\section{Timescale Separation and Relaxation Oscillations} \label{sec:oscillations}

The reward baseline $\bar{R}$ and synaptic weight $w$ evolve on timescales much slower than the remaining state variables. The baseline has $\tau_{\bar{R}} = 5$\,s, which is $250\times$ slower than the membrane time constant $\tau_m = 20$\,ms. The weight $w$ evolves according to $\dot{w} = M \cdot E$, where both the modulation signal $M = R - \bar{R}$ and the eligibility trace $E$ (scaled by learning rates $\eta_\pm \sim 10^{-4}$) are small in magnitude, yielding an effective timescale for $w$ that is comparable to or slower than $\tau_{\bar{R}}$.

\subsection{Spike-count quantization}

The postsynaptic neuron fires a discrete number of spikes $N(w)$ per presynaptic inter-spike interval $\Delta t_{\mathrm{pre}} = 1/r_{\mathrm{pre}}$. For a given weight $w$, the peak membrane drive is $R_m\, w\, I_s$, and the number of threshold crossings within each input cycle is a piecewise-constant, right-continuous function of $w$:
\[
    N(w) = \#\bigl\{t \in [t_{\mathrm{pre}}^k,\; t_{\mathrm{pre}}^{k+1}) : V(t) \uparrow \theta \bigr\}.
\]
At critical weights $w_c^{(n)}$ where $N$ increments from $n$ to $n+1$, the smoothed firing-rate trace $r_{\mathrm{post}}$ undergoes a step change of approximately $1/\tau_r$, and consequently the reward $R = -(r_{\mathrm{post}} - \alpha\, r_{\mathrm{pre}})^2$ experiences a discontinuous shift.

\subsection{Mechanism of the limit cycle}

Near a critical weight $w_c^{(n)}$, the system exhibits a relaxation oscillation with the following phases:

\paragraph{Slow recovery.} With $w < w_c^{(n)}$, the neuron fires $N = n$ spikes per input cycle. The reward $R$ and baseline $\bar{R}$ are approximately matched, so $M \approx 0$ and $\dot{w} \approx 0$. However, small residual $M > 0$ (since $\bar{R}$ lags $R$) causes $w$ to drift upward on the timescale $\mathcal{O}(\tau_{\bar{R}} / (\eta\, |E|))$.

\paragraph{Threshold crossing.} When $w$ crosses $w_c^{(n)}$, an additional postsynaptic spike fires. The discrete increment $r_{\mathrm{post}} \leftarrow r_{\mathrm{post}} + 1$ causes $R$ to jump sharply negative. Since $\bar{R}$ has not yet tracked this change ($\tau_{\bar{R}}$ is large), the modulation signal
\[
    M = R - \bar{R} = -\bigl(r_{\mathrm{post}} - \alpha\, r_{\mathrm{pre}}\bigr)^2 - \bar{R}
\]
becomes strongly negative. Combined with a nonzero eligibility trace $E$, this drives $\dot{w} = M \cdot E \ll 0$, rapidly decreasing $w$.

\paragraph{Subcritical reset.} With $w$ pushed below $w_c^{(n)}$, the extra spike no longer fires and $N$ returns to $n$. The reward $R$ recovers, and $\bar{R}$ begins to track back toward the (less negative) steady-state $R$. Both $\bar{R}$ and $w$ slowly recover, and the cycle repeats.

\subsection{Non-convergence}

The system cannot reach a stable fixed point at $w = w_c^{(n)}$ because the postsynaptic firing rate $r_{\mathrm{post}}(w)$ is discontinuous in $w$ at the critical values $w_c^{(n)}$, while the learning rule $\dot{w} = M \cdot E$ is continuous. The weight is attracted toward $w_c^{(n)}$ from below (where $M > 0$ drives $w$ up) and repelled from above (where the extra spike makes $M \ll 0$ and drives $w$ down). This constitutes a \emph{grazing bifurcation}: the smooth flow interacts with a discontinuous reset map at the spike threshold, generically producing a stable limit cycle rather than a fixed point.

\newpage
\section*{Equation Reference}

\paragraph*{Inter-spike dynamics --- input subsystem:}
\begin{align}
    \frac{dI_s}{dt} &= -\frac{I_s}{\tau_s}, \label{eq:ref_Is} \\[6pt]
    \frac{dx_{\mathrm{pre}}}{dt} &= -\frac{x_{\mathrm{pre}}}{\tau_+}, \label{eq:ref_xpre} \\[6pt]
    \frac{dr_{\mathrm{pre}}}{dt} &= -\frac{r_{\mathrm{pre}}}{\tau_r}. \label{eq:ref_rpre}
\end{align}

\paragraph*{Inter-spike dynamics --- postsynaptic:}
\begin{align}
    \tau_m \frac{dV}{dt} &= -(V - E_L) + R_m\, w\, I_s(t), \label{eq:ref_V} \\[6pt]
    \frac{dy_{\mathrm{post}}}{dt} &= -\frac{y_{\mathrm{post}}}{\tau_-}, \label{eq:ref_ypost} \\[6pt]
    \frac{dE}{dt} &= -\frac{E}{\tau_e}, \label{eq:ref_E} \\[6pt]
    \frac{dr_{\mathrm{post}}}{dt} &= -\frac{r_{\mathrm{post}}}{\tau_r}, \label{eq:ref_rpost} \\[6pt]
    \tau_{\bar{R}} \frac{d\bar{R}}{dt} &= -\bar{R} - \bigl(r_{\mathrm{post}} - \alpha\, r_{\mathrm{pre}}\bigr)^2, \label{eq:ref_Rbar} \\[6pt]
    \frac{dw}{dt} &= \Bigl[-\bigl(r_{\mathrm{post}} - \alpha\, r_{\mathrm{pre}}\bigr)^2 - \bar{R}\Bigr]\, E. \label{eq:ref_w}
\end{align}

\paragraph*{Pre-synaptic spike event at $t_{\mathrm{pre}}^k$:}
\begin{align}
    I_s &\;\leftarrow\; I_s + 1, \label{eq:ref_pre_Is} \\[6pt]
    x_{\mathrm{pre}} &\;\leftarrow\; x_{\mathrm{pre}} + 1, \label{eq:ref_pre_xpre} \\[6pt]
    r_{\mathrm{pre}} &\;\leftarrow\; r_{\mathrm{pre}} + 1, \label{eq:ref_pre_rpre}
\end{align}
\begin{align}
    E \;\leftarrow\; E - \eta_-\, w\, y_{\mathrm{post}}(t_{\mathrm{pre}}^{k-}). \label{eq:ref_pre_E}
\end{align}

\paragraph*{Post-synaptic spike event at $t_{\mathrm{post}}^k$:}
\begin{align}
    V &\;\leftarrow\; V_{\mathrm{reset}}, \label{eq:ref_post_V} \\[6pt]
    y_{\mathrm{post}} &\;\leftarrow\; y_{\mathrm{post}} + 1, \label{eq:ref_post_ypost} \\[6pt]
    r_{\mathrm{post}} &\;\leftarrow\; r_{\mathrm{post}} + 1, \label{eq:ref_post_rpost}
\end{align}
\begin{align}
    E \;\leftarrow\; E + \eta_+\,(w_{\mathrm{max}} - w)\, x_{\mathrm{pre}}(t_{\mathrm{post}}^{k-}). \label{eq:ref_post_E}
\end{align}

\section*{Model Parameters}
\label{sec:parameters}

\begin{table}[htbp]
\centering
\small
\label{tab:params}
\renewcommand{\arraystretch}{1.2}
\begin{tabular}{@{}lll>{\RaggedRight\arraybackslash}p{4.8cm}@{}}
\toprule
\textbf{Parameter} & \textbf{Description} & \textbf{Typical Value} & \textbf{Source} \\
\midrule
\multicolumn{4}{@{}l}{\textit{Membrane dynamics (leaky integrate-and-fire)}} \\[2pt]
$\tau_m$               & Membrane time constant          & 20\,ms                  &
  \cite{McCormick1985,kochMethodsNeuronalModeling1998,dayanTheoreticalNeuroscienceComputational2001} \\
$E_L$                  & Leak (resting) potential        & $-70$\,mV               &
  \cite{McCormick1985,Destexhe2001} \\
$R_m$                  & Membrane resistance             & $100$\,M$\Omega$        &
  \cite{kochMethodsNeuronalModeling1998,dayanTheoreticalNeuroscienceComputational2001} \\
$\theta$               & Spike threshold                 & $-55$\,mV               &
  \cite{Henze2000,Destexhe2001} \\
$V_{\mathrm{reset}}$   & Reset potential after spike     & $-65$\,mV               &
  \cite{Destexhe2001,gerstnerNeuronalDynamicsSingle2014} \\
$\tau_{\mathrm{ref}}$  & Absolute refractory period      & 2\,ms                   &
  \cite{Henze2000,gerstnerNeuronalDynamicsSingle2014} \\[6pt]
\multicolumn{4}{@{}l}{\textit{Synaptic current}} \\[2pt]
$\tau_s$               & Synaptic current decay (AMPA)   & 5\,ms                   &
  \cite{Destexhe1994,Roth2009,dayanTheoreticalNeuroscienceComputational2001} \\[6pt]
\multicolumn{4}{@{}l}{\textit{STDP trace time constants}} \\[2pt]
$\tau_+$               & Pre-synaptic LTP trace          & 20\,ms                  &
  \cite{biSynapticModificationsCultured1998,Song2000} \\
$\tau_-$               & Post-synaptic LTD trace         & 20\,ms                  &
  \cite{biSynapticModificationsCultured1998,Song2000} \\[6pt]
\multicolumn{4}{@{}l}{\textit{Firing-rate trace}} \\[2pt]
$\tau_r$               & Low-pass firing-rate trace      & 200\,ms                 &
  \cite{gerstnerNeuronalDynamicsSingle2014,izhikevichDynamicalSystemsNeuroscience2006} \\[6pt]
\multicolumn{4}{@{}l}{\textit{Eligibility trace}} \\[2pt]
$\tau_e$               & Eligibility trace decay         & 1\,s                    &
  \cite{gerstnerEligibilityTracesPlasticity2018,fremauxNeuromodulatedSpikeTimingDependentPlasticity2016} \\[6pt]
\multicolumn{4}{@{}l}{\textit{Reward baseline}} \\[2pt]
$\tau_{\bar{R}}$       & Running reward average          & 1--10\,s                &
  \cite{gerstnerEligibilityTracesPlasticity2018,fremauxNeuromodulatedSpikeTimingDependentPlasticity2016,Sutton1998} \\
$\alpha$               & Target firing-rate ratio        & 0.5                     &
  --- \\[6pt]
\multicolumn{4}{@{}l}{\textit{Plasticity learning rates}} \\[2pt]
$\eta_+$               & LTP learning rate               & $0.01$                  &
  \cite{Song2000,Morrison2008} \\
$\eta_-$               & LTD learning rate               & $0.0105$                &
  \cite{Song2000,Morrison2008} \\[6pt]
\multicolumn{4}{@{}l}{\textit{Synaptic weight bounds}} \\[2pt]
$w_{\mathrm{max}}$     & Maximum synaptic weight         & 1.0 (normalized)        &
  \cite{Song2000,Morrison2008} \\
$w(0)$                 & Initial synaptic weight         & $0.5\,w_{\mathrm{max}}$ &
  \cite{Song2000,abbottBuildingFunctionalNetworks2016} \\
\bottomrule
\end{tabular}
\end{table}

% \printbibliography

\end{document}